% ============================================================
% SUPERADO (03-ago-2026, misma sesion): la decision final de revista
% paso a Bioengineering (MDPI) -- ver
% docs/manuscrito/plantilla_overleaf_Bioengineering_MDPI.tex, esa es la
% plantilla vigente. Este archivo queda como referencia/alternativa, no
% se borra, por si se retoma POI mas adelante (ver
% docs/planificacion/revistas_candidatas_Q2.md).
% ============================================================
% ESQUELETO DE MANUSCRITO — listo para copiar a un proyecto nuevo de Overleaf
% Revista objetivo (version previa, ya no vigente): Prosthetics and Orthotics International (POI)
% Creado: 03-ago-2026
%
% ADVERTENCIAS IMPORTANTES — leer antes de usar:
%
% 1. NO se pudo verificar directamente la plantilla oficial de POI.
%    journals.sagepub.com bloquea el acceso automatizado (HTTP 403 en
%    author-instructions/poi y en home/poi). La busqueda web ademas dio
%    resultados CONTRADICTORIOS sobre el editor actual de la revista
%    (SAGE vs. Wolters Kluwer aparecen mezclados en distintas fuentes,
%    probablemente por listados de indexacion/Ovid desactualizados) —
%    el propio equipo ya habia verificado el 31-jul-2026 que
%    https://journals.sagepub.com/home/poi es la pagina oficial vigente
%    (ver docs/planificacion/revistas_candidatas_Q2.md), asi que este
%    esqueleto asume SAGE, pero CONFIRMAR ABRIENDO LA PAGINA EN EL
%    NAVEGADOR antes de enviar nada — el bloqueo es solo para fetch
%    automatizado, no para un navegador normal.
%
% 2. Este esqueleto usa la clase generica de SAGE para Overleaf
%    (sagej.cls, "A demonstration of the LaTeX2e class file for SAGE
%    Publications" — plantilla verificada en el propio Overleaf, no es
%    especifica de POI). Muchas revistas clinicas de SAGE (como POI)
%    piden el manuscrito en Word para el envio inicial y solo formatean
%    a la plantilla final tras aceptacion — verificar esto directamente
%    antes de asumir que un PDF compilado con sagej.cls es aceptado como
%    envio inicial.
%
% 3. Estilo de referencias sin confirmar: se dejo \bibliographystyle{sagev}
%    (SAGE-Vancouver, numerico, tipico de revistas clinicas/salud) como
%    mejor suposicion — CONFIRMAR contra la guia real de autores antes
%    de la version final. La alternativa es sagehvar (Harvard, autor-año).
%
% 4. Contenido: todo lo que ya esta redactado viene de
%    docs/manuscrito/metodos_introduccion_borrador.md (revisado
%    03-ago-2026, ya sin Kinovea/IMU de Alessandro/cita al paper de
%    conferencia). Los tramos [PENDIENTE: ...] no se completaron con
%    supuestos — mismo criterio que el resto del proyecto.
% ============================================================

% Opcion de clase sagev=Vancouver (numerico) -- mejor suposicion para
% revista clinica, SIN CONFIRMAR contra la guia real de POI (ver
% advertencia 3 arriba). Si POI usa Harvard, cambiar sagev por sageh y
% \bibliographystyle{sagev} por \bibliographystyle{sageh}.
\documentclass[times,sagev]{sagej}

\usepackage{amsmath}
\usepackage{graphicx}
\usepackage{booktabs}

\begin{document}

% ------------------------------------------------------------
% TITULO — tentativo, ver docs/planificacion/propuesta_articulo_Q2.md secc. 2
% ------------------------------------------------------------
\title{Multi-Subject Functional Validation of a 3-DOF Gait Simulator for
Transtibial Prosthesis Testing Using Inertial Motion Capture}

\runninghead{[PENDIENTE: apellidos de autores]}

% [PENDIENTE: nombres, afiliaciones y orden de autoria -- decidir con CRediT statement, ver checklist en propuesta_articulo_Q2.md secc. 7]
\author{Autor Uno\affilnum{1}, Autor Dos\affilnum{1}, Autor Tres\affilnum{2}, Autor Cuatro\affilnum{3}}

\affiliation{\affilnum{1}Departamento de Ingenier\'ia, Pontificia Universidad Cat\'olica del Per\'u (PUCP), Lima, Per\'u\\
\affilnum{2}[PENDIENTE: afiliaci\'on]\\
\affilnum{3}[PENDIENTE: afiliaci\'on]}

\corrauth{[PENDIENTE: autor de correspondencia], Laboratorio LIBRA,
Pontificia Universidad Cat\'olica del Per\'u, Lima, Per\'u.}
\email{[PENDIENTE: correo]}

\begin{abstract}
% [PENDIENTE: abstract completo — confirmar si POI exige formato
% estructurado (Background/Objective/Methods/Results/Conclusion) o
% abierto antes de redactar. Longitud tipica 200-250 palabras, confirmar
% limite exacto en la guia de autores.]
Background: [PENDIENTE].
Objective: [PENDIENTE].
Methods: [PENDIENTE].
Results: [PENDIENTE — depende de datos de sujetos nuevos y del piloto del IMU en la plataforma, ver CLAUDE.md].
Conclusion: [PENDIENTE].
\end{abstract}

\keywords{gait simulator, transtibial prosthesis, inertial motion capture, multi-subject validation, prosthetics testing}
% [PENDIENTE: confirmar numero maximo de keywords y si POI pide vocabulario controlado (MeSH)]

\maketitle

% ============================================================
\section{Introduction}
% ============================================================
% [PENDIENTE: cuerpo completo de la Introduccion — necesidad clinica de
% simuladores de marcha para evaluacion de protesis, con literatura
% externa (no autocita), antes del parrafo de cierre de abajo. Ver nota
% al inicio de metodos_introduccion_borrador.md sobre por que ya no se
% cita el paper de conferencia propio.]

Transtibial prosthesis testing increasingly relies on gait simulators to
provide controlled, repeatable loading conditions for functional and
mechanical evaluation without requiring continuous human-subject
involvement. A central open question for any such simulator is whether
trajectories programmed from a single reference subject generalize to
the gait patterns of individuals who did not contribute to that
programming --- a question that is rarely addressed with quantified,
multi-subject evidence in the literature on prosthesis gait simulators.
The present study addresses this gap directly: we assess whether a
3-degree-of-freedom (horizontal, vertical, sagittal) motor-driven gait
simulator reproduces gait kinematics captured from multiple subjects
independent of its calibration trajectory, using an inertial motion
capture system whose validity for lower-limb kinematics --- including
specifically in transtibial prosthesis users --- is already established
in the literature~\cite{PENDIENTE_iSen_transtibial}, and we quantify ---
rather than narrate --- the sources of vertical force overestimation
inherent to a motor-driven platform.

% ============================================================
\section{Methods}
% ============================================================

\subsection{System description}
% Condensado, autocontenido -- ya no remite al paper de conferencia (no citado en este articulo)

The gait simulator reproduces a pre-recorded gait trajectory across
three independent degrees of freedom --- horizontal displacement,
vertical displacement, and sagittal rotation --- each driven by its own
motor and transmission (chain-and-sprocket or direct drive, depending on
the axis). Trajectories are stored as CSV files, uploaded via a
Raspberry Pi, and executed open-loop by an ESP32 microcontroller; the
present study does not modify this operating principle --- no
closed-loop control or algorithmic trajectory generation was introduced
in this cycle (see Discussion, Future Work). % [PENDIENTE: frase sobre calibracion del actuador vertical, una vez resuelta la prueba de offset]

\subsection{Instrumentation}
% Revisado 03-ago-2026: un solo instrumento en todo el protocolo

Kinematic and kinetic data were collected using two instruments:

\begin{itemize}
\item \textbf{Inertial motion capture (STT-IWS/iSen).} Used throughout
the protocol to (i) capture new subjects' natural gait, (ii) re-capture
the original reference subject to regenerate the simulator's default
trajectory, and (iii) measure the simulator's own platform output
(sensor mounted directly on the platform). Segment/platform inclination
angles are derived from the sensor's onboard orientation estimate
relative to gravity, positive above horizontal and negative below, the
same sign convention used throughout this project. No in-house
cross-instrument validation was performed in this cycle; instrument
validity for lower-limb gait kinematics --- including specifically in
transtibial prosthesis users --- is supported by previously published
validation studies rather than re-derived~\cite{PENDIENTE_iSen_transtibial}.
% [PENDIENTE: confirmar que el sensor exporta orientacion cruda respecto
% a la gravedad y no solo el angulo articular calibrado del protocolo de
% marcha por defecto -- resultado del piloto en curso]
\item \textbf{Force platform (AMTI).} Used to record the vertical
ground reaction force (Fz) exerted by the simulator against a fixed
load cell, sampled at 1000\,Hz.
% [PENDIENTE: confirmar si se reportan las 6 columnas completas (Fx,Fy,Fz,Mx,My,Mz) o solo Fz]
\end{itemize}

\subsubsection{Vertical force overestimation: a three-stage, decoupled explanation}

Because the moving assembly is not a single rigid body --- a stationary
electronic control board is mounted on one side of the frame, while
three independently actuated axes (horizontal, vertical, sagittal) each
carry their own kinematic chain on the other --- vertical force
overestimation is decomposed into three independent, sequentially
corrected sources rather than a single inertial term:

\begin{enumerate}
\item A fixed \textbf{vertical datum offset}, calibrated once against
an independent static load reference (not against any subject
trajectory reported in Results) and held constant across all trials.
\item \textbf{Command-to-encoder tracking fidelity} per axis, isolating
mechanical/control error (chain or sprocket backlash, insufficient
motor holding torque under reaction load) from the remaining sources.
\item An \textbf{axis-wise inertial correction}, using the mass that
actually accelerates with each motor according to the real mounting
order of the kinematic chain --- not the total assembly mass ---
benchmarked against published prosthetic gait GRF data.
% [PENDIENTE: fuente de m_eje sin resolver -- CAD vs. renunciar al termino y declararlo limitacion]
\end{enumerate}

\subsection{Experimental protocol}
% [PENDIENTE: depende de si llego la aprobacion de etica y del numero final de sujetos -- no completar con supuestos, ver CLAUDE.md "Que falta decidir/hacer"]

\subsection{Statistical analysis}
% [PENDIENTE: prosa -- las formulas ya estan listas en
% docs/planificacion/plan_trabajo_5_semanas_articulo_Q2.md, seccion
% "Formulas estadisticas". Incluir: RMSEnorm, Pearson r (con la nota de
% cautela en curvas monotonas), % puntos dentro de +-1 SD, ICC(3,1)
% [Koo \& Li, 2016] especificando el modelo exacto, SPM1D no parametrico
% por permutacion [Nichols \& Holmes, 2002; Pataky et al., 2015]. NO
% mencionar Bland-Altman -- sin uso en este ciclo (ver
% docs/planificacion/propuesta_articulo_Q2.md secc. 4.1). Encuadrar el
% marco general de exactitud citando ISO 5725 (trueness/precision) --
% ver docs/literatura/normas_ISO_relevantes.md.]

% ============================================================
\section{Results}
% ============================================================
% Orden recomendado -- ver docs/planificacion/propuesta_articulo_Q2.md secc. 5

\subsection{Tracking fidelity per subject}
% [PENDIENTE: depende de datos de sujetos nuevos]

\subsection{Representativeness of the default trajectory}
% [PENDIENTE: depende de datos de sujetos nuevos]

\subsection{Repeatability and between-subject variability}
% [PENDIENTE: depende de datos de sujetos nuevos]

\subsection{Vertical force: raw, corrected, and literature benchmark}
% [PENDIENTE: depende de la calibracion del offset vertical, bloqueada por integracion Raspberry Pi-ESP32]

% ============================================================
\section{Discussion}
% ============================================================
% [PENDIENTE: interpretacion, limitaciones explicitas (tamano de
% muestra -- ver propuesta_articulo_Q2.md secc. 4.2 y 9), subseccion
% "Sources of error, decoupled", parrafo de trabajo futuro (lazo
% cerrado, generacion algoritmica de trayectorias, IMU de Alessandro y
% celda de carga para el segundo articulo).]

% ============================================================
\section{Conclusion}
% ============================================================
% [PENDIENTE]

% ============================================================
\section*{Declarations}
% ============================================================

\subsection*{Ethical approval}
% [PENDIENTE: numero de protocolo, depende de aprobacion del comite de etica]

\subsection*{Conflict of interest}
The authors declare no conflict of interest. % [PENDIENTE: confirmar]

\subsection*{Funding}
% [PENDIENTE]

\subsection*{Data and code availability}
% [PENDIENTE: decidir alcance -- repo completo publico vs. solo scripts de analisis con datos de ejemplo, ver propuesta_articulo_Q2.md secc. 7]

\subsection*{Author contributions (CRediT)}
% [PENDIENTE]

% ============================================================
\bibliographystyle{sagev} % [PENDIENTE: confirmar Vancouver vs. Harvard contra la guia real de POI]
\bibliography{referencias}
% Crear referencias.bib en el mismo proyecto de Overleaf. Entradas ya
% identificadas para incluir (ver docs/literatura/):
%   - validacion_instrumentos_IMU.md: validez iSen/STT-IWS (confirmar cita exacta antes de usar)
%   - validacion_instrumentos_IMU.md: IMU en usuarios de protesis transtibial (Sensors 2023, DOI 10.3390/s23031738)
%   - literatura_GRF_protesica.md: benchmark de vGRF en marcha protesica
%   - normas_ISO_relevantes.md: ISO 5725; Wu et al. 2002 (convencion ISB)
%   - Koo & Li (2016) -- ICC(3,1)
%   - Nichols & Holmes (2002); Pataky et al. (2015) -- SPM1D no parametrico
%   - Kadaba et al. (1989) -- CMC

\end{document}
